
% Encodage et polices
\usepackage[utf8]{inputenc}
\usepackage[T1]{fontenc}
\usepackage{lmodern, cmap}

% Bibliographie
\usepackage{natbib}

% Langue
\usepackage[french]{babel}

% Mathématiques
\usepackage{amsmath, amssymb, amsthm, mathtools, esint, dsfont, bm, mathrsfs, yfonts}

% Graphiques, couleurs et géométrie
\usepackage{graphicx, xcolor}
\usepackage{geometry}

% Tableaux et mise en page
\usepackage{array, multirow, diagbox}

% TikZ et PGFPlots
\usepackage{tikz}
\usetikzlibrary{decorations.pathmorphing, calc, shapes.geometric, positioning, arrows.meta, shapes, quotes}
\usepackage{pgfplots}
\pgfplotsset{compat=1.18}

% Figures, légendes et sous-figures
\usepackage{caption, subcaption}

% Algorithmes
\usepackage[ruled,noend]{algorithm2e}

% Divers utilitaires
\usepackage{multicol, paralist, microtype, url, framed}
\usepackage[framemethod=tikz]{mdframed} % utilisation unique avec l'option
\usepackage{ulem} % pour le soulignement


% Hyperliens et signets PDF (à charger en dernier)
\usepackage{hyperref}
\usepackage{bookmark}




\SetKwComment{Comment}{/* }{ */}
\DeclareMathOperator{\ones}{ones}
\DeclareMathOperator{\diag}{diag}
\DeclareMathOperator{\som}{sum}
\DeclareMathOperator{\tr}{Tr}



%========================================================================%
%                       Comptages et Décorations                         %
%========================================================================%




% Compteurs pour les environnements mathématiques
\newcounter{maths}  % Pour théorèmes, propositions, lemmes, corollaires
\newcounter{defi}   % Pour définitions


% Décoration ondulée
\newcommand\wavydecor{%
  \draw[%
    decoration={coil, aspect=0, segment length=5pt, amplitude=0.8pt},%
    decorate, line width=1.5pt, black%
  ] (O|-P) -- (O);
}

% Décoration linéaire
\newcommand\linedecor{%
  \draw[line width=0.8pt, black] (O|-P) -- (O);
}

% Décoration double barre
\newcommand\doublebardecor{%
  \draw[line width=0.8pt, black] (O|-P) -- (O);%
  \draw[line width=0.8pt, black] ([xshift=3pt]O|-P) -- ([xshift=3pt]O);
}


% Style théorème + variantes
\mdfdefinestyle{thmstyle}{%
  nobreak=false,
  hidealllines=true,%
  innerleftmargin=10pt, innerrightmargin=10pt,%
  innertopmargin=0pt, innerbottommargin=0pt,%
  leftmargin=10pt,%
  skipabove=.2\baselineskip, skipbelow=.5\baselineskip,%
  singleextra={\doublebardecor}, firstextra={\doublebardecor},%
  secondextra={\doublebardecor}, middleextra={\doublebardecor}%
}

% Style preuve
\mdfdefinestyle{proofstyle}{%
  nobreak=false,
  hidealllines=true,%
  innerleftmargin=10pt, innerrightmargin=10pt,%
  innertopmargin=0pt, innerbottommargin=0pt,%
  leftmargin=10pt,%
  skipabove=.2\baselineskip, skipbelow=.5\baselineskip,%
  font=\itshape,%
  singleextra={\linedecor}, firstextra={\linedecor},%
  secondextra={\linedecor}, middleextra={\linedecor}%
}


% Style définition
\mdfdefinestyle{defstyle}{%
  nobreak=false,
  hidealllines=true,%
  innerleftmargin=10pt, innerrightmargin=10pt,%
  innertopmargin=0pt, innerbottommargin=0pt,%
  leftmargin=10pt,%
  skipabove=.2\baselineskip, skipbelow=.5\baselineskip,%
  singleextra={\wavydecor}, firstextra={\wavydecor},%
  secondextra={\wavydecor}, middleextra={\wavydecor}%
}



% Théorème, Proposition, Lemme, Corollaire 

\newenvironment{theo}[2][]{%
    \vspace{10pt}\par\noindent\refstepcounter{maths}%
    \ifx\relax#2\relax\else\label{#2}\fi%
    \textsc{Théorème~\themaths}%
    \ifx\relax#1\relax\else\ (\textsc{#1})\fi :%
    \begin{mdframed}[style=thmstyle]%
}{%
    \end{mdframed}%
}

\newenvironment{prop}[2][]{%
    \vspace{10pt}\par\noindent\refstepcounter{maths}%
    \ifx\relax#2\relax\else\label{#2}\fi%
    \textsc{Proposition~\themaths}%
    \ifx\relax#1\relax\else\ (\textsc{#1})\fi :%
    \begin{mdframed}[style=thmstyle]%
}{%
    \end{mdframed}%
}

\newenvironment{lem}[2][]{%
    \par\noindent\refstepcounter{maths}%
    \ifx\relax#2\relax\else\label{#2}\fi%
    \textsc{Lemme~\themaths}%
    \ifx\relax#1\relax\else\ (\textsc{#1})\fi :%
    \begin{mdframed}[style=thmstyle]%
}{%
    \end{mdframed}%
}

\newenvironment{exo}[2][]{%
    \vspace{10pt}\par\noindent\refstepcounter{maths}%
    \ifx\relax#2\relax\else\label{#2}\fi%
    \textsc{Exercice~\themaths}%
    \ifx\relax#1\relax\else\ (\textsc{#1})\fi :%
    \begin{mdframed}[style=thmstyle]%
}{%
    \end{mdframed}%
}


\newenvironment{coro}[2][]{%
    \par\noindent\refstepcounter{maths}%
    \ifx\relax#2\relax\else\label{#2}\fi%
    \textsc{Corollaire~\themaths}%
    \ifx\relax#1\relax\else\ (\textsc{#1})\fi :%
    \begin{mdframed}[style=thmstyle]%
}{%
    \end{mdframed}%
}

% Preuve

\newenvironment{preuve}[1][]{%
    \noindent\textit{Preuve :} \ifx\relax#1\relax\else\ \textit{#1}\fi\par
    \begin{mdframed}[style=proofstyle]%
}{%
    \hfill$\square$\end{mdframed}%
}


% Environnement pour Définition

\newenvironment{defi}[2][]{%
    \vspace{10pt}\par\noindent\refstepcounter{defi}%
    \ifx\relax#2\relax\else\label{#2}\fi%
    \textsc{Définition~\thedefi}%
    \ifx\relax#1\relax\else\ (\textsc{#1})\fi :%
    \begin{mdframed}[style=defstyle]%
}{%
    \end{mdframed}%
}




% Nouvelles commandes et raccourcis

\renewcommand{\leq}{\leqslant}
\renewcommand{\geq}{\geqslant}
\newcommand{\ms}[1]{\mathscr{#1}}
\newcommand{\Phx}{\Phi_x}
\newcommand{\Phy}{\Phi_y}
\newcommand{\X}{\mathcal{X}}
\newcommand{\Y}{\mathcal{Y}}
\newcommand{\Z}{\mathcal{Z}}

% Commandes d'usage spécifique
\newcommand{\norm}[1]{\left\|#1\right\|}    % norme
\newcommand{\C}{\mathscr{C}}    % classe C
\renewcommand{\S}{\mathbb{S}}   % sphère
\newcommand{\R}{\mathbb{R}}     % Réels
\newcommand{\Ra}{\mathscr{R}}   % Radon
\renewcommand{\d}{\textnormal{d}}     % différentiel
\renewcommand{\P}{\mathscr{P}}  % polynôme
\renewcommand{\hat}{\widehat}   % chapeaux de bonne taille
\newcommand{\ch}{\wedge}
\newcommand{\cc}{{\ch \ch}}
\newcommand{\ce}{{\ch *}}
\newcommand{\ec}{{* \ch}}
\newcommand{\ee}{{* *}}
\renewcommand{\Pr}{\mathbb{P}}

\newcommand{\N}{\mathbb{N}}
\newcommand{\I}{\mathcal{I}}
\newcommand{\U}{\mathcal{U}}
\newcommand{\E}{\mathbb{E}}
\renewcommand{\H}{\mathscr{H}}
\newcommand{\G}{\mathscr{G}}
\newcommand{\fat}{f_\theta^a}
\newcommand{\fbt}{f_\theta^b}

\newcommand{\nhs}[1]{\left\|#1\right\|_\text{HS}}
\newcommand{\scal}[2]{\left\langle #1 \vert #2 \right\rangle}
\newcommand{\floor}[1]{\left\lfloor#1\right\rfloor}
\newcommand{\indep}{\perp\!\!\!\perp}

\everymath{\displaystyle}

%\usepackage{libertinus}